\section{Reproducibility}\label{sec:repro}

The harness, the adapters, the plans, the frozen operating points and \texttt{FREEZE.json} are in
the repository~\cite{sweetsearch_repo} under \texttt{eval/benchmarks/hnsw-rivals/paper} (frozen
commit \texttt{f57533b0}, tag \texttt{hnsw-rivals-freeze-2026-10-17-v2}). The release adds the
vectors of every set and model, the splits, the ground truth, every held-out result record with its
retained lists, and the per-job timing files. From these, one command rebuilds each table and
figure:

\begin{small}\begin{verbatim}
python analyze.py --data DATA --results RESULTS --shortlists shortlists \
    --builds BUILDS --types T --out R-T.md --json J-T.json  # per CPU T
python merge_results.py --parts . --out RESULTS-paper.md \
    --json results-paper.json
python mem_patch.py --results RESULTS --json results-paper.json
python draft/mktex.py results-paper.json --extra draft/extra.json
python draft/figures.py results-paper.json
\end{verbatim}\end{small}

\texttt{freeze\_check.py} confirms that the frozen files are unchanged, and
\texttt{check\_complete.py} confirms that every planned held-out job has a complete record
(\nm{hojobs} jobs on the three AWS CPUs). Adding a system takes one adapter with the common
interface (build, load, search, return a vector list with scores) and its build and search grids in
\texttt{grid.py}; the selection and analysis scripts need no change. \TODO{release kit: a small laptop example (gcsn, one model, two systems), its measured run time, and the archive DOI}
